\section{結言}
今回の観測窓は、Navier--Stokes問題への解という極めて大きな研究主張と、個人エージェント、画像生成、コーディング、サービング基盤の製品更新が同時に重なった。特にOpenAIの数学発表は、モデル性能だけでなく、多数エージェントの協調と形式検証を組み合わせる研究様式そのものを前面に出している。ただし、その数学的正しさや賞の認定は今後の独立検証に委ねられる。

一方、Muse、Mercury 2.5、Cohere megakernel、CursorへのSpark展開は、生成AI競争がモデル単体からエージェント製品、IDE、推論基盤へ広がっていることを示した。未確認の抗生物質説や個人実験も含め、同じ24時間の中に確定度の異なる情報が混在しており、Daily Xではその差を保ったまま「何が話題になったか」を記録することが重要な観測日だった。
